\chapter{Conclusion} \label{cap:conclusion}

Kloud Motors achieved the main goal of the project: building a \textbf{cloud-native online car dealer platform} around a realistic vehicle dataset. The final system is not only an API over static data, but a distributed application composed of a public gateway, internal gRPC services, PostgreSQL databases, Redis caching, Firebase authentication, WebSocket-based real-time features, Pub/Sub event distribution, Kubernetes deployment manifests, Terraform infrastructure, monitoring, backups, and CI/CD automation.

The \textbf{microservices architecture} proved useful for separating the main business domains of the platform. Listing, search, market price, geographical insights, users, sellers, chat, and auctions can evolve independently while still being exposed through one public REST and WebSocket gateway. This structure also made it possible to introduce cross-cutting concerns in a controlled way: the gateway centralises public authentication and HTTP error mapping, services keep their internal gRPC contracts, Redis reduces repeated read pressure, and circuit breakers and timeouts help the system behave better when dependencies are slow or unavailable.

The deployment work was an important part of the project. The application can run locally with Docker Compose, but it is also prepared for cloud execution on GKE through Kubernetes manifests and Kustomize. Terraform defines the supporting Google Cloud resources, including Cloud SQL, Artifact Registry, Pub/Sub, Cloud Storage, service accounts, networking, and the GKE cluster. The CI/CD workflows connect these pieces by running Go tests, executing Docker Compose integration tests, applying Terraform changes, building and pushing service images, and rolling the workloads out to GKE.

The Cloudy Day and stress-test phase exposed several practical issues that are common in distributed systems. These tests motivated improvements in gateway validation, gRPC timeouts, circuit breakers, database indexes, resource requests and limits, HPA configuration, and operational documentation. As a result, the final version is more robust than the initial implementation and better aligned with the non-functional requirements of scalability, resilience, observability, and maintainability.

There are still limitations. The GKE node pool and Cloud SQL instance are intentionally small to keep the assignment cost under control, so the current deployment should be seen as a \textit{realistic prototype} rather than a production-sized system. The HPA configuration prepares services for horizontal scaling, but a production deployment would also require node autoscaling, stronger database capacity planning, stricter secret management, TLS configuration, image tag immutability, and more complete end-to-end tests for authenticated and real-time workflows.

Overall, the project demonstrates the \textbf{full lifecycle of a cloud application}: data preparation, API design, service implementation, containerisation, cloud infrastructure provisioning, Kubernetes deployment, automated validation, monitoring, backup, and cost-aware operation. The result is a working platform that satisfies the main functional requirements while also exercising the operational practices expected from a modern cloud computing system.
